% page 14 du fac-simile (image p-013 du scan)
% bloc 160.0 x 242.0 mm ; marge gauche 8.0 mm
\pgc{-12.700mm}{0.000mm}{
\l{\cel{6}6}
\l{}
\l{}
\l{\sou{acidometro}.\cel{2}Instrumento per qua onu mezuras la koncentreso-}
\l{\cel{3}grado di acido. - EFIRS.}
\l{}
\l{\sou{aciono.}- (\sou{Financo})\cel{2}I. Parto di la kapitalo suskriptenda,}
\l{\cel{3}qua yurizas pri parto di la trezoro societala e di la}
\l{\cel{3}profito di societo. - II. La akto qua reprezentas singla}
\l{\cel{3}di ta parti. - DFIRS.}
\l{}
\l{\sou{acipitro}. (\sou{zool}.) Rapt-ucelo jornala, qua ne flugas en la}
\l{\cel{3}alta strati di atmosfero. - L.\cel{1}\sou{astur palumbarius}.}
\l{}
\l{\sou{acizo}.\cel{2}Imposto di ula vari (nutrivi, e c.) ye la enireyo di}
\l{\cel{3}urbo. - DEFIRS.}
\l{}
\l{\sou{adaptar}.\cel{2}(\sou{trans., ad}). Unionar (kozo) ad altra kozo qua kon-}
\l{\cel{3}venas a lu. - DEFIRS.}
\l{}
\l{\sou{adenino}.\cel{2}(\sou{kemio})\cel{2}C5H5N5.}
\l{}
\l{\sou{adenito}. (\sou{patol}.) Inflamuro di la glandi o ganglioni limfala.}
\l{}
\l{\sou{adepto}.- Persono quan onu iniciis pri ula arto o doktrino.}
\l{\cel{3}- DEFIRS.}
\l{}
\l{\sou{adequata}.- (\sou{filoz}.)Qua equivalas, esas adaptita, korespondas}
\l{\cel{3}perfekte a la objekto (aludante penso, cienco, e c.p - DEFS.}
\l{}
\l{\sou{adherar}.\cel{2}(\sou{netrans., ad})\cel{2}I. Quaze glutinesar an ulo.-}
\l{\cel{3}II. (\sou{metaf}.) Ligar su komplete ad opiniono, a partiso.-DEFIS}
\l{}
\l{\sou{adiar}. (\sou{trans}).I. Pronuncar la formulo di politeso kande onu}
\l{\cel{3}su separas (de ulu). - II. Pronuncar la formulo di renunco}
\l{\cel{3}od abandono. - DEFIS.}
\l{}
\l{\sou{adiabata}. (\sou{fiziko}) Quan kaloro povas permear.}
\l{}
\l{\sou{adicionar}\cel{2}(\sou{matem.})\cel{2}Adjuntar nombro ad un o plura nombri di}
\l{\cel{3}la sama naturo, por unionar li ad una. - DEFIS.}
\l{}
\l{\sou{adipocero}. (\sou{kemio)}.\cel{2}Nomo olima di tri bbstanci quin onu kredis}
\l{\cel{3}inter-identa, ma qui interdiferas esencale : spermaceto,}
\l{\cel{3}la grasajo di la kadavri, e kolesterino (C27 H45 OH).}
\l{}
\l{\sou{adjektivo}.\cel{2}(\sou{gram}.) Vorto-speco qua expresas la eso-maniero.}
\l{\cel{3}- DEFIS.}
\l{}
\l{\sou{adjudikar}.\cel{2}(\sou{trans., ad}). I. (\sou{Yurocienco)}\cel{2}Atribuar, per judi-}
\l{\cel{3}cio, a la autoritato o persono qui esas kompetenta.- II.}
\l{\cel{3}(\sou{lor vendo}) Atribuar a la persono qua ofras po lo la preco}
\l{\cel{3}maxima. - DEFIS.}
\l{}
\l{\sou{adjurar}. (\sou{trans., pri})\cel{2}Sumnar (ulu), ye la nomo di sakrajo}
\l{\cel{3}ke lu agez ulo. - EF}
}
